\documentclass[11pt]{article}
 
\usepackage[T1]{fontenc}
\usepackage[utf8]{inputenc}
\usepackage[margin=1in]{geometry}
\usepackage{amsmath,amssymb,amsthm,mathtools}
\usepackage{enumitem}
\usepackage{booktabs}
\usepackage{longtable}
\usepackage{array}
\usepackage[hidelinks]{hyperref}
 
\newtheorem{definition}{Definition}[section]
\newtheorem{lemma}[definition]{Lemma}
\newtheorem{proposition}[definition]{Proposition}
\newtheorem{theorem}[definition]{Theorem}
\newtheorem{remark}[definition]{Remark}
\newtheorem{convention}[definition]{Convention}
 
\newcommand{\R}{R}
\newcommand{\Sfam}{\mathcal S}
\newcommand{\Lean}[1]{\texttt{\begingroup\def\_{\textunderscore\discretionary{}{}{}}#1\endgroup}}
\newcommand{\Const}{\operatorname{Const}}
\newcommand{\Inj}{\operatorname{Inj}}
\newcommand{\lmin}{\operatorname{lmin}}
\newcommand{\Int}{\operatorname{Int}}
 
\title{Detailed mathematical write-up of the Lean formalisation\\
Lemma 1 of the Monotonicity Theorem: a constant or injective subinterval}
\author{}
\date{}
 
\begin{document}
\maketitle
 
\section{Purpose and high-level structure}
 
This document explains, in ordinary mathematical language, a Lean formalisation of \emph{Lemma~1}
from the proof of the Monotonicity Theorem for definable functions on an o-minimal structure. In
the source notes (\Lean{monotonicity.tex}), the statement is recorded as a lemma occurring inside
the proof of the Monotonicity Theorem, and is referred to in the running text simply as Lemma~1.
 
Informally, the argument splits into two cases according to whether some value of $f$ is attained
infinitely often on the interval $I$. If so, that infinite fibre already contains a subinterval on
which $f$ is constant. Otherwise every fibre of $f$ on $I$ is finite, which forces the image
$f(I)$ to be infinite; choosing, for each $y$ in an interval inside $f(I)$, the \emph{least}
preimage $g(y)\in I$, produces a one-sided inverse of $f$ that is injective by construction. The
image of $g$ must again contain a subinterval of $I$, and on this last subinterval $f$ is
necessarily injective.
 
The Lean file discussed here (\Lean{kovacs\_benedek\_noel.lean}) formalises exactly this two-case
argument. Unlike the companion development described in \emph{Detailed mathematical write-up of
the Lean formalisation: O-minimal structures, definable maps, images, and local properties}, the
present file does not build or import a full \Lean{OMinimalStructure} object indexed by all
arities $n$. Instead, it isolates the two non-trivial mathematical facts that the argument for
Lemma~1 actually needs and records them as two explicit Lean axioms:
\begin{itemize}
  \item \Lean{o\_minimal\_contains\_interval}: every infinite subset of $R$ contains an infinite
  subinterval of $R$. This is precisely the one-dimensional content of o-minimality --- Definition
  5.1(xi) in the companion document, together with density of the order --- specialised to the
  case at hand.
  \item \Lean{infinite\_image\_of\_injOn}: the image of an infinite set under a function that is
  injective on that set is again infinite. This is an ordinary fact of cardinal arithmetic; it
  carries no o-minimality content whatsoever.
\end{itemize}
Sections~\ref{sec:order-and-functions}--\ref{sec:choice-function} introduce the order-theoretic
vocabulary and these two axioms, together with short proofs explaining \emph{why} each axiom is
true of an actual o-minimal structure. Section~\ref{sec:proof} then gives the proof of Lemma~1
itself, following the Lean tactic proof step by step.
 
The main result of the file is the following statement.
 
\begin{lemma}[Lemma~1: a constant or injective subinterval]\label{lem:main}
Let $R$ be a nonempty linearly ordered set, let $I\subseteq R$ be an infinite interval, and let
$f:R\to R$ be a definable function. Then there is a subinterval $J\subseteq I$ such that $f$ is
either constant on $J$ or injective on $J$.
\end{lemma}
 
This is the Lean theorem \Lean{lemma\_0\_5}.
 
\section{The ambient order, intervals, and definable functions}\label{sec:order-and-functions}
 
\subsection{Intervals}
 
Throughout, $R$ is a type equipped with a linear order $<$ (with associated $\le$), and all sets
are subsets of $R$. The Lean file works with the following order-theoretic notion of an interval.
 
\begin{definition}[Interval]\label{def:interval}
A subset $S\subseteq R$ is an \emph{interval}, written $\Int(S)$, if it is order-convex: for all
$x,y\in S$ and all $z\in R$,
\[
  x\le z\le y \implies z\in S.
\]
\end{definition}
 
In Lean this is the predicate \Lean{is\_interval}.
 
\begin{remark}
Definition~\ref{def:interval} is the usual order-theoretic notion of a convex subset and is
slightly more permissive than the inductively generated class \Lean{FiniteUnionOfPointsAndIntervals}
of the companion document (which singles out points and \emph{open} intervals as generators).
Every point, every open interval, and every closed or half-open interval satisfies
Definition~\ref{def:interval}. This extra generality is harmless here: the proof of Lemma~1 never
inspects the internal structure of an interval, only its convexity, so the weaker and more
flexible Definition~\ref{def:interval} is the natural notion to use in this file.
\end{remark}
 
\subsection{Definable functions}
 
\begin{definition}[Definable function, placeholder form]\label{def:definable-fn}
In this file, a function $f:R\to R$ is \emph{definable} if it satisfies the always-true predicate
$\Lean{DefinableFunction}(f) :\equiv \top$.
\end{definition}
 
\begin{remark}
In the full theory (Definition~6.1 of the companion document), definability of $f:A\to B$ is a
substantive condition: it requires $A$ and $B$ to be definable sets and the graph of $f$ to be a
definable subset of $R^{2}$. Because the present file does not instantiate an ambient
\Lean{OMinimalStructure}, the predicate \Lean{DefinableFunction} is left as the trivial placeholder
of Definition~\ref{def:definable-fn}. Its purpose is purely structural: it lets the statement of
Lemma~\ref{lem:main} keep the same shape, $f$ ranging over ``definable'' functions, that it would
have in the full theory, so that \Lean{DefinableFunction} can later be replaced by the genuine
notion of Definition~6.1 without changing the statement or proof of Lemma~\ref{lem:main} at all.
\end{remark}
 
\section{Constancy and injectivity on a subset}
 
\begin{definition}[Constant and injective on a set]
Let $f:R\to R$ and $J\subseteq R$. We write
\[
  \Const(f,J) :\iff \forall x,y\in J,\ f(x)=f(y),
\]
\[
  \Inj(f,J) :\iff \forall x,y\in J,\ \bigl(f(x)=f(y)\implies x=y\bigr).
\]
\end{definition}
 
In Lean these are \Lean{ConstantOn} and \Lean{InjectiveOn}. With this vocabulary,
Lemma~\ref{lem:main} asserts exactly that some subinterval $J\subseteq I$ satisfies
$\Const(f,J)\vee\Inj(f,J)$.
 
\section{The o-minimal input: infinite sets contain an interval}
 
\begin{definition}[Infinite-set axiom]\label{def:omin-axiom}
The file postulates, as a Lean \Lean{axiom}, that for every infinite $S\subseteq R$ there is a
subinterval $J\subseteq S$ that is itself infinite:
\[
  S\subseteq R,\ S\text{ infinite}\ \Longrightarrow\ \exists J\subseteq S,\ \Int(J)\wedge J\text{
  infinite}.
\]
\end{definition}
 
This is \Lean{o\_minimal\_contains\_interval}. Although it is introduced as an axiom in this
self-contained file, it is not an independent assumption about $R$: it is a theorem of any actual
o-minimal structure, and the next lemma records its proof.
 
\begin{lemma}\label{lem:omin-derivation}
Let $M=(R,<,(\Sfam_n)_{n})$ be an o-minimal structure on a dense linear order without endpoints
$R$ (Definitions~3.1 and~5.1 of the companion document). If $S\in\Sfam_1$ is infinite, then $S$
contains an infinite open interval.
\end{lemma}
 
\begin{proof}
By the one-dimensional o-minimality axiom (Definition~5.1(xi)), $S$ is a finite union of points
and open intervals,
\[
  S = \{a_1,\dots,a_k\}\ \cup\ (b_1,c_1)\ \cup\ \cdots\ \cup\ (b_m,c_m).
\]
The set of points $\{a_1,\dots,a_k\}$ is finite. If there were no interval pieces ($m=0$), $S$
itself would be finite, contradicting the hypothesis. Hence $m\ge1$, and $S$ contains some open
interval $(b_j,c_j)$ with $b_j<c_j$. By density of the order (Definition~3.1(iv)), between any two
elements of $(b_j,c_j)$ there is a further element of $(b_j,c_j)$; iterating this produces
infinitely many distinct elements, so $(b_j,c_j)$ is infinite. Since an open interval is convex,
$J:=(b_j,c_j)$ satisfies $\Int(J)$, $J\subseteq S$, and $J$ is infinite, as required.
\end{proof}
 
\begin{remark}
Lemma~\ref{lem:omin-derivation} is exactly Definition~\ref{def:omin-axiom}, i.e.\
\Lean{o\_minimal\_contains\_interval}, specialised at $n=1$ and restated without reference to the
family $(\Sfam_n)_n$. Recording it as an axiom in the present file is a deliberate simplification:
it avoids re-deriving, inside this file, the inductive structure of points and open intervals
underlying \Lean{Finite\discretionary{}{}{}Union\discretionary{}{}{}Of\discretionary{}{}{}Points\discretionary{}{}{}AndIntervals},
together with the density argument, while
remaining faithful to the underlying mathematics, since the displayed implication is a genuine
theorem of every o-minimal structure.
\end{remark}
 
\section{A cardinality fact: injective images of infinite sets}
 
\begin{definition}[Injective-image axiom]
The file also postulates, as a Lean \Lean{axiom}, that for $s\subseteq R$ and $f:R\to R$,
\[
  s\text{ infinite},\ f\text{ injective on } s\ \Longrightarrow\ f(s)\text{ infinite}.
\]
\end{definition}
 
This is \Lean{infinite\_image\_of\_injOn}. In contrast with
Definition~\ref{def:omin-axiom}, this fact carries no o-minimality content: it is a statement of
ordinary cardinal arithmetic, true for an arbitrary function between arbitrary sets.
 
\begin{lemma}\label{lem:inj-image}
Let $s$ be an infinite set and let $f$ be injective on $s$. Then $f(s)$ is infinite.
\end{lemma}
 
\begin{proof}
We argue by contraposition. Suppose $f(s)$ is finite. The restriction of $f$ to $s$, viewed as a
map $s\to f(s)$, is injective by hypothesis and surjective onto $f(s)$ by definition of the image,
hence a bijection of $s$ onto $f(s)$. An infinite set cannot be put in bijection with a finite one,
so $s$ would have to be finite as well. This contradicts the hypothesis that $s$ is infinite.
Hence $f(s)$ is infinite.
\end{proof}
 
\begin{remark}
Because Lemma~\ref{lem:inj-image} does not depend on o-minimality, it could equally well be proved
inside the Lean file itself rather than postulated as an axiom (for instance by formalising the
contrapositive argument above). It is recorded here as an axiom purely for economy of
presentation; doing so does not weaken the formalisation of Lemma~\ref{lem:main}, since
Lemma~\ref{lem:inj-image} is unconditionally true.
\end{remark}
 
\section{A choice function for least elements}\label{sec:choice-function}
 
The informal proof of Lemma~1 defines an ``inverse'' $g$ by $g(y):=\min\{x\in I: f(x)=y\}$. To
formalise this without first proving that the relevant set always has a minimum, the Lean file
uses Hilbert's choice operator.
 
\begin{definition}[Least element, by choice]
For $S\subseteq R$, define
\[
  \lmin(S) := \varepsilon x.\ \bigl(x\in S \wedge \forall y\in S,\ x\le y\bigr),
\]
where $\varepsilon$ is Hilbert's choice operator: $\varepsilon x.\,P(x)$ denotes an arbitrarily
chosen element satisfying $P$ if one exists, and an arbitrary element of $R$ otherwise.
\end{definition}
 
This is the noncomputable Lean definition \Lean{myInf}, built from \Lean{Classical.epsilon}.
Nothing is asserted yet about $\lmin(S)$ in general; the next lemma isolates exactly the
circumstances under which $\lmin(S)$ behaves as a genuine minimum.
 
\begin{lemma}[Specification of $\lmin$]\label{lem:myinf-spec}
If $S$ is finite and nonempty, then $\lmin(S)\in S$ and $\lmin(S)\le z$ for every $z\in S$; that
is, $\lmin(S)$ is the minimum element of $S$.
\end{lemma}
 
\begin{proof}
Since $R$ is linearly ordered and $S$ is finite and nonempty, $S$ has a minimum element $x_0$: view
$S$ as a finite set (a \Lean{Finset}) and let $x_0$ be its minimum, which exists by finiteness and
is characterised by $x_0\in S$ and $x_0\le z$ for all $z\in S$. Thus the predicate
$P(x):\equiv x\in S\wedge\forall y\in S,\ x\le y$ is satisfied by $x_0$, so $\exists x,\ P(x)$ holds.
By the defining specification of Hilbert's choice operator, whenever $\exists x,\ P(x)$ holds, the
chosen element $\varepsilon x.\,P(x)$ itself satisfies $P$. Hence $P(\lmin(S))$ holds, which is the
statement of the lemma.
\end{proof}
 
\begin{raggedright}
This is the Lean lemma \Lean{myInf\_spec}. In Lean the minimum $x_0$ is produced by
\Lean{Set.Finite.toFinset} together with \Lean{Finset.min\textquotesingle}. The final step,
\Lean{Classical.epsilon\_spec}, then transfers this minimality to $\lmin$.
\end{raggedright}
 
\section{Proof of Lemma~1}\label{sec:proof}
 
We now prove Lemma~\ref{lem:main}, following the case split used in the Lean proof of
\Lean{lemma\_0\_5}. Fix $R$, $I$, $f$ as in the statement, with $I$ an infinite interval.
 
\subsection{Splitting on the fibres of $f$ over $I$}
 
The proof distinguishes two cases, according to whether
\[
  \exists y\in R,\ f^{-1}(y)\cap I \text{ is infinite}
\]
holds or fails. In Lean this is the case split \Lean{by\_cases h\_inf\_preimage}.
 
\subsection{Case 1: some fibre is infinite}
 
Suppose $f^{-1}(y)\cap I$ is infinite for some $y\in R$. By
Definition~\ref{def:omin-axiom}/\Lean{o\_minimal\_contains\_interval} applied to the infinite set
$f^{-1}(y)\cap I$, there is a subinterval
\[
  J\subseteq f^{-1}(y)\cap I
\]
with $\Int(J)$ (the auxiliary infiniteness of $J$ furnished by the axiom is not needed here and is
discarded; in Lean this is \Lean{h\_contains\_interval}). Since $J\subseteq I$, $J$ is a candidate
subinterval of $I$. We show $\Const(f,J)$. For $x,z\in J$ we have $x,z\in f^{-1}(y)$, i.e.
$f(x)=y=f(z)$; hence $f(x)=f(z)$. As $x,z\in J$ were arbitrary, $\Const(f,J)$ holds, and the
disjunction $\Const(f,J)\vee\Inj(f,J)$ in Lemma~\ref{lem:main} is witnessed by its left disjunct.
 
\subsection{Case 2: every fibre is finite}
 
Now suppose every fibre of $f$ on $I$ is finite:
\[
  \forall y\in R,\ f^{-1}(y)\cap I \text{ is finite}.
\]
(In Lean, this is obtained from the negation of the Case~1 hypothesis by the tactic \Lean{push
Not}, using the fact that, by definition, a set is infinite exactly when it is not finite, so the
negated hypothesis simplifies directly to finiteness of every fibre; this hypothesis is recorded
as \Lean{h\_inf\_preimage} for the remainder of the proof.)
 
\subsubsection{The image $f(I)$ is infinite}
 
\begin{lemma}\label{lem:fI-infinite}
Under the standing hypotheses of Case~2, $f(I)$ is infinite.
\end{lemma}
 
\begin{proof}
Suppose, for contradiction, that $f(I)$ is finite. Every $x\in I$ lies in the fibre
$f^{-1}(f(x))\cap I$, so
\[
  I \ \subseteq\ \bigcup_{y\in f(I)} \bigl(f^{-1}(y)\cap I\bigr).
\]
The index set $f(I)$ is finite by assumption, and each set $f^{-1}(y)\cap I$ in the union is
finite by the standing hypothesis of Case~2. A finite union of finite sets is finite, so the
right-hand side is finite, and hence $I$, being a subset of a finite set, is finite. This
contradicts the hypothesis that $I$ is infinite.
\end{proof}
 
This is \Lean{h\_fI\_inf}; in Lean the contradiction is reached directly, by assuming a proof
\Lean{h\_fin\_img} of finiteness of $f(I)$, exhibiting $I$ as a subset of the finite union
\Lean{Set.Finite.biUnion h\_fin\_img (fun y \_ => h\_inf\_preimage y)} via
\Lean{Set.mem\_biUnion}, and finally applying the hypothesis \Lean{hI\_inf} (that $I$ is infinite,
i.e.\ that $I$ is \emph{not} finite) to the resulting proof that $I$ is finite.
 
By Definition~\ref{def:omin-axiom}/\Lean{o\_minimal\_contains\_interval} applied to the infinite
set $f(I)$, we obtain a subinterval
\[
  J_{\mathrm{img}}\subseteq f(I), \qquad \Int(J_{\mathrm{img}}), \qquad J_{\mathrm{img}}\text{
  infinite}.
\]
This is \Lean{h\_fI\_contains\_interval}; unlike in Case~1, the infiniteness of $J_{\mathrm{img}}$
furnished by the axiom is retained here, since it will be needed below.
 
\subsubsection{A one-sided inverse $g$ on $J_{\mathrm{img}}$}
 
Define
\[
  g:R\to R,\qquad g(y) := \lmin\bigl(\{x\in I : f(x)=y\}\bigr).
\]
This is the Lean local definition \Lean{let g := fun y => myInf \{x \(\in\) I | f x = y\}}.
 
\begin{lemma}\label{lem:g-properties}
For every $y\in J_{\mathrm{img}}$,
\[
  g(y)\in I \qquad\text{and}\qquad f(g(y)) = y.
\]
\end{lemma}
 
\begin{proof}
Fix $y\in J_{\mathrm{img}}\subseteq f(I)$, so there is $x\in I$ with $f(x)=y$; thus the set
$S_y:=\{x\in I: f(x)=y\}$ is nonempty. By the standing hypothesis of Case~2, $S_y$ is a subset of
the finite fibre $f^{-1}(y)\cap I$ (indeed $S_y = f^{-1}(y)\cap I$), so $S_y$ is finite. By
Lemma~\ref{lem:myinf-spec}, $\lmin(S_y) = g(y)$ satisfies $g(y)\in S_y$, i.e. $g(y)\in I$ and
$f(g(y))=y$.
\end{proof}
 
This is \Lean{h\_g\_in\_I} and \Lean{h\_fg\_eq}; the finiteness of each $S_y$ is recorded once, for
all $y$ simultaneously, as \Lean{h\_finite\_preimage}.
 
\subsubsection{$g$ is injective, and its image contains an interval}
 
\begin{lemma}\label{lem:g-injective}
$g$ is injective on $J_{\mathrm{img}}$.
\end{lemma}
 
\begin{proof}
Let $y_1,y_2\in J_{\mathrm{img}}$ with $g(y_1)=g(y_2)$. By Lemma~\ref{lem:g-properties},
\[
  y_1 = f(g(y_1)) = f(g(y_2)) = y_2.
\]
\end{proof}
 
This is \Lean{h\_g\_inj}. Combining Lemma~\ref{lem:g-injective} with the infiniteness of
$J_{\mathrm{img}}$ established above and Lemma~\ref{lem:inj-image}, i.e.\
\Lean{infinite\_image\_of\_injOn}, the image $g(J_{\mathrm{img}})$ is infinite; this is
\Lean{h\_gJ\_inf}. Applying
Definition~\ref{def:omin-axiom}/\Lean{o\_minimal\_contains\_interval} once more, this time to the
infinite set $g(J_{\mathrm{img}})$, yields a subinterval
\[
  J \subseteq g(J_{\mathrm{img}}), \qquad \Int(J)
\]
(again discarding the auxiliary infiniteness of $J$, which is not needed further). This is
\Lean{h\_gJ\_contains\_interval}.
 
By Lemma~\ref{lem:g-properties}, every element of $g(J_{\mathrm{img}})$ lies in $I$, so
$J\subseteq g(J_{\mathrm{img}})\subseteq I$. This inclusion is \Lean{hJ\_sub\_I}, and it makes $J$
a legitimate candidate subinterval of $I$ for Lemma~\ref{lem:main}.
 
\subsubsection{Conclusion: $f$ is injective on $J$}
 
It remains to show $\Inj(f,J)$. Let $x_1,x_2\in J$ with $f(x_1)=f(x_2)$. Since
$J\subseteq g(J_{\mathrm{img}})$, there are $y_1,y_2\in J_{\mathrm{img}}$ with
\[
  g(y_1)=x_1, \qquad g(y_2)=x_2.
\]
Substituting, the hypothesis $f(x_1)=f(x_2)$ becomes
\[
  f(g(y_1)) = f(g(y_2)),
\]
and the desired conclusion $x_1=x_2$ becomes, after the same substitution, the equivalent goal
\[
  g(y_1) = g(y_2).
\]
By Lemma~\ref{lem:g-properties}, $f(g(y_1))=y_1$ and $f(g(y_2))=y_2$; combined with
$f(g(y_1))=f(g(y_2))$ from above, this gives $y_1=y_2$. Substituting back, $g(y_1)=g(y_2)$, which
is exactly the goal. Hence $\Inj(f,J)$ holds, and the disjunction $\Const(f,J)\vee\Inj(f,J)$ is
witnessed by its right disjunct. This completes Case~2, and with it the proof of Lemma~\ref{lem:main}.
 
In Lean, the two substitutions described above are performed by rewriting along
\Lean{hy1\_eq : g y1 = x1} and \Lean{hy2\_eq : g y2 = x2} (in the hypothesis \Lean{h\_eq} and in
the goal), the equality $y_1=y_2$ is recorded as \Lean{h\_y\_eq}, obtained by rewriting along
\Lean{h\_fg\_eq y1 hy1\_in} and \Lean{h\_fg\_eq y2 hy2\_in}, and the proof finishes with
\Lean{rw [h\_y\_eq]}.
 
\subsection{A remark on the hypotheses \texorpdfstring{$\Int(I)$}{Int(I)} and definability of $f$}
 
\begin{remark}
The Lean signature of \Lean{lemma\_0\_5} records two hypotheses, named \Lean{\_hI : is\_interval I}
and \Lean{\_hf\_def : DefinableFunction f}, whose leading underscore is the Lean convention for
``present for the shape of the statement, not used by name in the proof.'' Indeed, neither the
convexity of $I$ nor the (placeholder) definability of $f$ is used anywhere in the argument above:
only the infiniteness of $I$, recorded as \Lean{hI\_inf}, is needed. This is not an oversight; it
mirrors the informal proof, which likewise never uses convexity of $I$ beyond the fact that a
subset of an interval that is itself convex is again an interval --- a fact that is used only to
justify calling the sets $J$ produced along the way ``subintervals,'' not to derive any further
property of $f$. The hypotheses are kept in the statement of Lemma~\ref{lem:main} so that its
signature matches the way $I$ and $f$ are introduced in the surrounding proof of the Monotonicity
Theorem, where $I$ is always an interval and $f$ is always definable.
\end{remark}
 
\section{Dictionary between the Lean file and this document}
 
\begin{longtable}{>{\raggedright\arraybackslash}p{0.46\textwidth}>{\raggedright\arraybackslash}p{0.44\textwidth}}
\toprule
Lean name & Mathematical meaning \\
\midrule
\endhead
\Lean{is\_interval} & The convexity predicate $\Int(S)$ of Definition~\ref{def:interval}. \\
\Lean{DefinableFunction} & Placeholder predicate, identically true (Definition~\ref{def:definable-fn}). \\
\Lean{ConstantOn} & The predicate $\Const(f,J)$. \\
\Lean{InjectiveOn} & The predicate $\Inj(f,J)$. \\
\Lean{o\_minimal\_contains\_interval} & Axiom: every infinite subset of $R$ contains an infinite subinterval (Definition~\ref{def:omin-axiom}; proved from o-minimality in Lemma~\ref{lem:omin-derivation}). \\
\Lean{infinite\_image\_of\_injOn} & Axiom: the injective image of an infinite set is infinite (Lemma~\ref{lem:inj-image}). \\
\Lean{myInf} & The choice function $\lmin(S)=\varepsilon x.\,(x\in S\wedge\forall y\in S, x\le y)$. \\
\Lean{myInf\_spec} & Specification of $\lmin$ on finite nonempty sets (Lemma~\ref{lem:myinf-spec}). \\
\Lean{lemma\_0\_5} & Lemma~1 itself (Lemma~\ref{lem:main}). \\
\Lean{h\_inf\_preimage} & The case-split hypothesis on fibres of $f$ over $I$ (Section~\ref{sec:proof}.1). \\
\Lean{h\_contains\_interval} & The subinterval of an infinite fibre, in Case~1. \\
\Lean{h\_fI\_inf} & Infiniteness of $f(I)$ in Case~2 (Lemma~\ref{lem:fI-infinite}). \\
\Lean{h\_fI\_contains\_interval} & The subinterval $J_{\mathrm{img}}\subseteq f(I)$, with its infiniteness. \\
\Lean{g} & The one-sided inverse $g(y)=\lmin\{x\in I: f(x)=y\}$. \\
\Lean{h\_finite\_preimage} & Finiteness of $\{x\in I: f(x)=y\}$ for every $y$, from the Case~2 hypothesis. \\
\Lean{h\_g\_in\_I}, \Lean{h\_fg\_eq} & The two parts of Lemma~\ref{lem:g-properties}: $g(y)\in I$ and $f(g(y))=y$. \\
\Lean{h\_g\_inj} & Injectivity of $g$ on $J_{\mathrm{img}}$ (Lemma~\ref{lem:g-injective}). \\
\Lean{h\_gJ\_inf} & Infiniteness of $g(J_{\mathrm{img}})$. \\
\Lean{h\_gJ\_contains\_interval} & The final subinterval $J\subseteq g(J_{\mathrm{img}})$. \\
\Lean{hJ\_sub\_I} & The inclusion $J\subseteq I$. \\
\Lean{h\_y\_eq} & The equality $y_1=y_2$ used to conclude $g(y_1)=g(y_2)$ at the end of Case~2. \\
\bottomrule
\end{longtable}
 
\section{Summary in one line}
 
The mathematical content of the Lean file is exactly the case split
\[
  \bigl(\exists y,\ f^{-1}(y)\cap I \text{ infinite}\bigr)
  \ \vee\
  \bigl(\forall y,\ f^{-1}(y)\cap I \text{ finite}\bigr),
\]
together with the one-dimensional o-minimality fact that every infinite subset of $R$ contains an
infinite subinterval, applied first to a single infinite fibre, and, in the remaining case, to the
infinite image $f(I)$ and again to the infinite image of the resulting least-preimage function
$g$; this assembles, in either branch, a subinterval $J\subseteq I$ on which $f$ is constant or
injective.
 
\end{document}